% !PRINTED 318
\weberpage{318}

In order to find an upper limit for the absolute value of the imaginary
roots which remains valid also in the case of complex coefficients, we
consider the function
\[
f(x)=x^n+a_1x^{n-1}+a_2x^{n-2}+\cdots+a_n,
\]
in which the coefficients $a_1, a_2\ldots a_n$ are real or imaginary.
We denote the absolute value of a complex quantity $a$, as before, by
$|a|$, and then, if we assume that the absolute value of $x$ is greater
than 1, we have
\[
\left| a_1+\frac{a_2}{x}+\cdots+\frac{a_n}{x^{n-1}} \right| <
|a_1|+|a_2|+\cdots+|a_n|=\alpha.
\]
If we therefore take $x$ so that its absolute value is greater than
each of the two numbers 1 and $\alpha$, the expression
\[
x^{1-n}f(x)=x+a_1+\frac{a_2}{x}+\cdots+\frac{a_n}{x^{n-1}}
\]
cannot vanish, since the absolute value of $x$ is greater than the
absolute value of the sum of all the remaining terms of the right-hand
side.

If, then, according as $\alpha$ is greater or less than 1, we draw a
circle with the radius $\alpha$ or 1 about the origin as centre,
certainly no roots of $f(x)$ lie outside this circle any longer.

\begin{center}
\S.\ 105.

\gesperrt{Estimation of the imaginary roots.}
\end{center}

We have seen in \S.\,96 how, from the characteristic of a certain
system of curves, we can determine the exact number of the complex
roots of an equation that lie in the interior of a closed curve. This
criterion too can be simplified if one no longer wishes to find the
exact number of the roots, but only an upper limit for this number,
which in many cases also suffices for the exact determination.

Let, then, as in \S.\,96
\weq{1}{F(x)=\varphi(x,y)+i\psi(x,y)}
be a real or complex function of the complex argument
%CONT>>
